\section{Contest quick reference}
\textbf{Before coding:} identify constraints, graph direction, indexing,
overflow bounds, duplicates, and whether empty input is possible. Sketch
the invariant and a small example. Estimate memory before allocating.

\textbf{Before submitting:} run the sample; try $n=1$, empty ranges, all equal,
negative values, disconnected components, repeated edges, and maximum values.
Reset state for each case. Check output precision and trailing case labels.

\textbf{Local compile:}
\begin{lstlisting}[language=bash]
g++ -std=c++17 -O2 -Wall -Wextra main.cpp -o main
g++ -std=c++17 -g -O1 -fsanitize=address,undefined \
  -fno-omit-frame-pointer main.cpp -o main
\end{lstlisting}
\textbf{Stress testing:} write a slow oracle for small inputs; generate random
cases with a fixed seed; compare results and save the first failing input.

\subsection{Vim setup and commands}
If you use Vim, type these settings into \path{~/.vimrc} and open
\texttt{main.cpp}. They use two-space indentation, keep namespace contents
unindented like the templates, and use stock Vim features;
no plugins are needed. In an already open Vim, load them with
\texttt{:source \textasciitilde/.vimrc}. The source ZIP includes
\path{codebook/contest.vimrc} for practice.
\lstinputlisting[language={}]{codebook/contest.vimrc}
\begin{tabular}{@{}ll@{}}
\texttt{Esc} & Return to normal mode \\
\texttt{i / o} & Insert / new line below \\
\texttt{:w / :wq} & Save / save and quit \\
\texttt{u / Ctrl-r} & Undo / redo \\
\texttt{yy / p / dd} & Copy / paste / delete line \\
\texttt{/text, n / N} & Search, next / previous match \\
\texttt{* / \%} & Search current word / matching brace \\
\texttt{gg=G} & Reindent the whole file \\
\texttt{:noh} & Clear search highlighting \\
\texttt{Ctrl-n} (insert mode) & Complete a word from buffers
\end{tabular}

\textbf{Compile and run inside Vim:} save first with \texttt{:w}, then use
the commands below (for \texttt{main.cpp} in the current directory).
\begin{lstlisting}[language={}]
:!g++ -std=c++17 -O2 -Wall -Wextra main.cpp -o main
:!./main < input.txt
\end{lstlisting}
If terminal pasting adds unwanted indentation, use \texttt{:set paste}
before pasting and \texttt{:set nopaste} afterward.

\subsection{Choose an algorithm}
\begin{tabular}{@{}ll@{}}
Unweighted shortest path & BFS \\
Weights 0 or 1 & 0-1 BFS \\
Nonnegative weights & Dijkstra \\
Negative weights & Bellman-Ford \\
All pairs, small graph & Floyd-Warshall \\
Static range min/max/gcd & Sparse table \\
Point add / prefix sum & Fenwick \\
Range update / query & Lazy segment tree \\
Tree paths / subtrees & HLD + segment tree \\
Boolean binary clauses & 2-SAT \\
Bipartite pairs & Hopcroft-Karp \\
Capacity / disjoint paths & Flow \\
Many pattern counts & Aho-Corasick \\
Palindrome substrings & Manacher \\
Polynomial multiplication & NTT \\
Linear DP transitions & Li Chao
\end{tabular}

\subsection{C++ edge cases}
\texttt{long long} holds about $9.22\cdot10^{18}$. Widen before multiplying:
\texttt{1LL * a * b} or \texttt{(\_\_int128) a * b}. Signed overflow is undefined.
Use \texttt{numeric\_limits<T>::lowest()} for a maximum-query identity.
Use integer literals for large sentinels. Never add an infinity sentinel.

\texttt{\_\_builtin\_clz(0)} and \texttt{\_\_builtin\_ctz(0)} are undefined.
Use \texttt{1LL << k} with a valid shift count; do not shift signed 1 into its
sign bit. \texttt{vector.size()} is unsigned. Negative modulo is negative:
normalize with \texttt{x \%= m; if (x < 0) \{ x += m; \}}.

\texttt{lower\_bound}: first element $\ge x$.
\texttt{upper\_bound}: first element $>x$.
Sort before using them. Comparators must be strict; do not use $\le$.
Erase/insert may invalidate iterators. Recursive DFS on a long chain may
exhaust the stack; use an explicit stack where needed.

\subsection{Compression and binary search}
\begin{lstlisting}
auto xs = a;
sort(xs.begin(), xs.end());
xs.erase(unique(xs.begin(), xs.end()), xs.end());
int id = (int) (lower_bound(xs.begin(), xs.end(), x) - xs.begin());
// First true, if pred is false...false,true...true.
// Answer n means no true element; search domain [0,n).
long long lo = 0, hi = n;
while (lo < hi) {
  long long mid = lo + (hi - lo) / 2;
  if (pred(mid)) {
    hi = mid;
  } else {
    lo = mid + 1;
  }
}
\end{lstlisting}
Compression preserves order, not distances. For geometric lengths use
the original coordinate differences.
